\section{Expanded reasoning: recognition and writer identity}
\label{sec:expand-recognise}
\subsection{Input dimensions, scale and information loss}
The inspected input is one grey channel on a $64\times640$ canvas.
For source height $h$ and width $w$, isotropic fit uses
$a=\min(64/h,640/w)$, producing rounded dimensions near $(ah,aw)$.
White padding fills the rest. Isotropic means the same scale in both axes;
it preserves aspect ratio before rounding. Stretching directly to the canvas
would distort letters. The report's $156\times2170$ example is width-limited:
$a=640/2170\simeq0.29493$, so the content height is about 46 pixels.
That is why a large nominal canvas height does not guarantee a tall letter
image. A very long line can be compressed until dots or gaps disappear.

The exact canvas size is a project architecture choice and a checkpoint
compatibility constraint, not a theorem. Increasing width retains more detail
but increases feature/sequence work; increasing height can preserve vertical
shape but changes activation memory. A trained fixed-shape model cannot be
assumed accurate under a new preparation convention. Variable-width batches
are an alternative, with padding/length masking costs. Pixel range, ink polarity,
placement and stroke normalisation must match training, not merely dimensions.

\subsection{What local feature map $Z$ actually means}
For input $X_{b,c,y,x}$, weights $W_{o,c,i,j}$ and bias $b_o$, a stride-one
layer computes
\[
Z_{b,o,y,x}=b_o+\sum_{c=1}^{C_{\rm in}}\sum_{i,j}
 W_{o,c,i,j}X_{b,c,y+i-p,x+j-p}.
\]
$b$ indexes batch examples, $c$ input channels, $o$ output channels,
$(i,j)$ the kernel and $p$ zero padding. Each $Z_o$ is a learned response
map, not a labelled character map. The operation used in machine learning
is cross-correlation (the kernel is not flipped), conventionally called
convolution~\cite{torchconv}. Sharing the same weights at every position
expresses the assumption that a useful local stroke pattern can appear anywhere.

Constructed $2\times2$ example: a patch
$\left[\begin{smallmatrix}1&2\\3&4\end{smallmatrix}\right]$ and kernel
$\left[\begin{smallmatrix}1&0\\0&-1\end{smallmatrix}\right]$ with bias 0.5
give $1-4+0.5=-2.5$. ReLU would then give zero. The kernel values in a
real network are learned, so visual intuition about an edge detector does
not establish what any particular saved channel detects.

\paragraph{Shape and parameter arithmetic.}
For kernel $k$, padding $p$, stride $s$ and dilation $d$,
$H_{\rm out}=\lfloor(H+2p-d(k-1)-1)/s+1\rfloor$; apply the same formula
to width. With $k=3,p=1,s=d=1$, size is unchanged. A convolution with bias
has $C_{\rm out}(C_{\rm in}k^2+1)$ parameters; $1\to32$ channels costs
$32(9+1)=320$, whereas $32\to32$ costs $32(288+1)=9248$.
More channels increase representation capacity and both parameter and
activation costs. Channel counts 32, 64 and 128 are implementation design
choices; powers of two are convenient, not evidence of an optimum.

The reported blocks use two, four and six convolutions, with two $2\times2$
spatial pools. Output shapes are $32\times32\times320$,
$64\times16\times160$ and $128\times16\times160$.
Taking a maximum over height yields $128\times1\times160$; transposing
gives 160 ordered feature vectors, each of length 128. A frame is therefore
a feature column, not a presegmented character or a video frame.
The CNN/recurrent/transcription combination follows the CRNN family
of Shi et al.~\cite{shi2017}. Its use does not justify these exact widths.

\paragraph{Receptive field.}
Track receptive-field span $r$ and source-pixel jump $j$. Starting at
$(r,j)=(1,1)$, a layer of kernel $k$ and stride $s$ gives
$r'=r+(k-1)j$, $j'=sj$. For the reported stack: two $3$-kernels give
$r=5$, first pool gives $(6,2)$; four more convolutions give 22, second
pool gives $(24,4)$; six more give $r=72,j=4$. Thus each final spatial
position can depend on about a $72$-pixel input span, subject to padding.
The LSTM adds longer horizontal context. A 4-pixel frame stride does not
mean a character boundary is localised to four source-image pixels.

\subsection{Batch normalisation, ReLU, sigmoid and max pooling}
For one channel containing $m=BHW$ values, training normalisation uses
$\mu=m^{-1}\sum_i z_i$, $v=m^{-1}\sum_i(z_i-\mu)^2$ and
$\hat z_i=(z_i-\mu)/\sqrt{v+\epsilon}$; the learned affine output is
$u_i=\gamma\hat z_i+\beta$~\cite{ioffe2015}.
$\gamma$ changes scale, $\beta$ restores a useful offset, and $\epsilon$
prevents unstable division for almost-constant input. For $(1,3)$,
$\mu=2,v=1$, so ignoring tiny $\epsilon$, normalised values are $(-1,1)$.
$\gamma=2,\beta=0.5$ produces $(-1.5,2.5)$.
Inference uses saved running statistics, not statistics of the current single
crop. A running-update coefficient controls how quickly history is forgotten;
it is not the optimiser learning rate. Libraries can use different variance
conventions for the running estimate~\cite{torchbn}.

ReLU is $\max(0,u)$, derivative 1 for $u>0$ and 0 for $u<0$, with an
implementation convention at zero. It makes the network nonlinear while
keeping the positive branch unsaturated~\cite{glorotrect}. Without nonlinear
operations, compositions of affine maps reduce to a single affine map.
This observation does not mean a network containing max pooling would become
linear merely by removing ReLU: max pooling is itself nonlinear.

Sigmoid $\sigma(z)=1/(1+e^{-z})$ has derivative $\sigma(z)(1-\sigma(z))$.
It maps real scores into $(0,1)$; $\sigma(0)=0.5$ and
$\sigma(\log3)=0.75$. It is used as a gate because multiplication can
attenuate or pass information. A gate is not automatically a calibrated
probability of a real-world event. Saturation near 0 or 1 yields small
derivatives. Tanh maps into $(-1,1)$ and has derivative $1-\tanh^2z$.

Max pooling selects the largest value in each window. For
$\left[\begin{smallmatrix}1&4\\3&2\end{smallmatrix}\right]$ it returns 4.
The backward pass routes the upstream derivative to the selected maximum,
using a fixed tie convention. Pooling reduces work and offers limited local
shift tolerance, while discarding exact position. Average pooling retains
average strength; strided convolution learns a reduction. Neither is
automatically preferable for small character details.

\subsection{LSTM gates, bidirectionality and the second recurrent layer}
In the gate convention used by the inspected implementation,
\[
\begin{bmatrix}a_i\\a_f\\a_g\\a_o\end{bmatrix}
=W_xx_t+W_hh_{t-1}+b,\quad
i_t=\sigma(a_i),\ f_t=\sigma(a_f),\ g_t=\tanh(a_g),\ o_t=\sigma(a_o),
\]
\[
c_t=f_t\odot c_{t-1}+i_t\odot g_t,\qquad
h_t=o_t\odot\tanh(c_t).
\]
$\odot$ means elementwise multiplication. The cell $c_t$ carries a memory
value; $h_t$ exposes a gated version. Weights are shared across frames and
do not get ``transferred'' from frame to frame. Input $x_t$, hidden state
and cell state do. The original LSTM motivation was learning long temporal
dependencies~\cite{hochreiter1997}; practical gate conventions are documented
in the framework reference~\cite{torchlstm}.

Constructed scalar cell: $c_{t-1}=0.8,f_t=0.9,i_t=0.2,g_t=0.5,o_t=0.7$
gives $c_t=0.82$ and $h_t\simeq0.472$. The forget gate retains 90\% of
the previous cell along this direct path; other dependencies also affect the
total gradient. A BiLSTM processes the same feature sequence left-to-right
and right-to-left and concatenates hidden outputs. With 256 units in each
direction it outputs 512 values per frame. The second layer therefore takes
512 inputs, whereas the first takes 128. Offline line images supply both
directions; a live causal system would not have future context immediately.

For input width $D$ and hidden width $h$, one direction has
$4h(D+h)$ weights and either $4h$ combined biases or $8h$ separately stored
input/recurrent biases. Counting the latter twice as two \emph{effective}
offsets would misdescribe the affine sum. Increasing hidden width has a
quadratic recurrent cost. The selected 256 and two layers are architecture
settings, not uniquely implied by CTC or handwriting.

\subsection{CTC: paths, blank, collapse and a complete small example}
Let $p_t(k)$ be the softmax probability of symbol $k$ at frame $t$.
For a path $\pi$ of $T$ symbols, its conditional score is
$\prod_t p_t(\pi_t)$. Collapse $B$ first merges consecutive identical
symbols, then removes blank. Thus $(a,a,\varnothing,b,b)$ becomes $ab$,
whereas $(a,\varnothing,a)$ becomes $aa$. Blank is a special alignment
symbol, not the space character. A literal space is an ordinary label.
CTC sums the scores of \emph{all} paths collapsing to the target and takes
negative logarithm~\cite{graves2006}.

Constructed two-frame example, target $a$: at frame one probabilities for
$(a,\varnothing)$ are $(0.6,0.4)$ and at frame two $(0.7,0.3)$.
The valid paths are $(a,a),(a,\varnothing),(\varnothing,a)$.
Their total is $0.42+0.18+0.28=0.88$, so loss is $-\log0.88\simeq0.1278$.
The all-blank path scores 0.12 and produces the empty string. Computing
only the best path would give 0.42 and a different training loss.

For target length $U$ with $r$ adjacent repeated pairs, at least $U+r$
frames are needed: each repeated adjacent label needs an intervening blank.
For ``book'', $U=4,r=1$, so five frames suffice in principle. This condition
is necessary, not evidence the network will learn good alignment.

\paragraph{Dynamic summation, not a separate dynamic collapse algorithm.}
Expand target labels to $\ell'=(\varnothing,\ell_1,\varnothing,\ldots,
\ell_U,\varnothing)$, giving $2U+1$ states. A transition can stay, advance
one state, or skip one blank when the new label is nonblank and differs
from the label two states back. Forward probability $\alpha_{t,s}$ sums
valid prefixes ending at $s$, multiplying by the current emission.
Initialise only the first blank and first label; terminate by summing the
last label and last blank states. Log-domain arithmetic replaces sums by
log-sum-exp to avoid underflow. For values $a_i$, use
$m+\log\sum_i e^{a_i-m}$ with $m=\max_i a_i$. Collapse itself is just
the deterministic repeat/blank removal rule~\cite{ctcdistill}.

Greedy inference chooses $\arg\max_k p_t(k)$ independently at each frame,
then collapses. It maximises a single path, not necessarily the total score
of a collapsed transcript. Beam search considers several prefixes and can
sum path scores; a language model can add linguistic assumptions. The
reported reader uses greedy decoding without a language model.

\subsection{Forced alignment and reading the CTC plot}
Given a known transcript, forced alignment uses the same allowed state
transitions but maximises the sum of log probabilities. Write
$V_{t,s}=\log p_t(\ell'_s)+\max_{r\in\mathcal P(s)}V_{t-1,r}$.
Store predecessors and backtrack. This Viterbi-style path is a timing/location
explanation conditional on the supplied transcript, not a fresh proof that
the transcript is correct. If a model is weak, a legal alignment can still
put cuts inside letters. Forward--backward soft occupancy and hard best-path
intervals answer different questions.

A frame edge $t$ maps approximately to canvas coordinate $640t/T$, then
to source $x=(640t/T)/a$ for left-aligned fit. With $T=160,a=0.5$, edge
10 maps to source $x=80$. Padding offsets, rounding and endpoint conventions
must be included when different preparation is used. CNN receptive fields
and CTC's flexible blank placement mean this is an approximate geometric
correspondence, not a character bounding-box measurement.

In the report's CTC plot, horizontal position is feature frame, each row a
class, and the shown trace can be either a greedy class or a transcript-state
alignment. Repeated $g$ labels need a blank-separated path. Blank carrying
about 46\% of average probability mass is compatible with correct decoding:
many frames are gaps or redundant representations. Probability mass is not
the fraction of blank image pixels. Compare labels, axes and the two paths
before interpreting the picture as segmentation evidence.

\subsection{Writer head, $K$ classes, dropout and stroke normalisation}
The writer head averages the 160 feature columns, applies a 96-unit hidden
layer with ReLU and a final $K$-logit layer. $K$ is the number of writers in
the loaded mapping, not the text alphabet size. The report evaluates eleven;
current path lists can contain thirteen. Checkpoint class order determines
which output index means which writer. A softmax over a closed set has no
automatic ``unknown writer'' class. Pooling loses order but summarises
distributed stroke features; it does not prove writer identity is independent
of text content.

With dropout probability $p=0.2$, training uses
$\tilde h=m h/(1-p)$ for Bernoulli keep mask $m$ of mean 0.8.
Therefore $E[\tilde h]=h$. Inference removes the random mask. Dropout is
published regularisation~\cite{srivastava2014}; 0.2 is the project's chosen
strength, not a universal optimum. For a 96-unit layer the expected number
kept is 76.8, not exactly 77 on every pass.

For the shape classifier, binary ink is thinned and square-dilated to controlled
width: $w=\operatorname{clip}(\operatorname{round}(\rho h_x),2,9)$,
$k=2\lfloor w/2\rfloor+1$, then white-background image
$I_s=255[1-\delta_k(\operatorname{thin}(M))]$. At $h_x=40,\rho=0.17$,
$w=7,k=7$; at $h_x=20$, $w=3,k=3$. If $w=2$, the odd side is 3.
Thus the effective width is quantised, not exactly $\rho h_x$. Square
dilation uses a Chebyshev neighbourhood; a disk is a different metric.
Reducing pen-width cues can make geometry more comparable, but preprocessing
can also erase genuine writer cues and amplify segmentation errors.

The report records widths 0.06, 0.09, 0.11 and 0.17, best epochs 38, 43,
39 and 43, and displayed accuracies 98.06\%, 97.42\%, 98.06\% and
99.35\%. These support ``best observed among these four settings'' for
that selection experiment. The surrounding prose calls it page identification
while Table 4 says per-line accuracy: the unit is unresolved here and must
not silently be chosen. It is not an independent post-selection qualification
score. The archived run settings and held-out unit would be needed to resolve
the discrepancy. No stroke-width sweep was rerun.
